\documentclass[12pt]{article}
\usepackage[margin=1in]{geometry}
\usepackage{amsmath}
\usepackage{fancyhdr}
\fancypagestyle{plain}{%
  \fancyhf{}
  \fancyhead[L]{POLISCI 450A}
  \fancyhead[R]{Fall 2026}
  %\fancyfoot[C]{\thepage}
}

\usepackage{titlesec}
\titleformat{\section}{\large\bfseries}{\thesection}{1em}{}
\usepackage{titling}
\setlength{\droptitle}{-3em}   % move title up

\usepackage{xcolor}
\newcommand{\red}{\color{red}}

\title{\Large Lecture 2-2: Interval Estimation}
\author{}
\date{}
\begin{document}
\maketitle
\vspace{-6em}

\noindent \textbf{Goal:} Obtain confidence intervals using normal or $t$ approximations.\medskip

\begin{itemize}\setlength{\itemsep}{1em}
\item Interval estimation
\begin{itemize}
\item Point vs. interval estimates
\item Interpreting confidence intervals
\item Coverage rate
\end{itemize}
\item Normal mean with known variance
\begin{itemize}
\item Standardization with known $\sigma$
\item $z$-based confidence intervals
\item Choosing confidence levels
\end{itemize}
\item Normal mean with unknown variance
\begin{itemize}
\item Estimating $\sigma^2$ via sample variance
\item Student's $t$ distribution
\item $t$-based confidence intervals
\end{itemize}
\item Mean with unknown distribution
\begin{itemize}
\item Central limit theorem justification
\item Asymptotic normal intervals
\item Prefer $t$ for small samples
\end{itemize}
\item Practical guidance
\begin{itemize}
\item When to use $z$ vs. $t$
\item Conservative choices by sample size
\end{itemize}
\end{itemize}


% End outline
\clearpage


\section*{Worksheet}



\begin{enumerate}\itemsep1.2em

\item (Two variance estimators) Define
$S_{0n}^2=\frac{1}{n}\sum_{i=1}^n (X_i-\overline{X}_n)^2$ and
$S_{1n}^2=\frac{1}{n-1}\sum_{i=1}^n (X_i-\overline{X}_n)^2$. Show:
\begin{itemize}
    \item[i)] $E[S_{0n}^2]=\frac{n-1}{n}\sigma^2$ and $E[S_{1n}^2]=\sigma^2$.
    \item[ii)] Both estimators are consistent. (Hint: write $S_{0n}^2=\frac{1}{n}\sum_{i=1}^n(X_i-\mu)^2-(\overline{X}_n-\mu)^2$ and apply the WLLN to each term.)
\end{itemize}
\vspace{18em}


\item ($t$ statistic) With unknown $\sigma^2$ and $X_i\sim N(\mu,\sigma^2)$, prove
$T_n=\frac{\overline{X}_n-\mu}{S_{1n}/\sqrt{n}} \sim t_{n-1}.$
(Hint: If $Z \sim N(0,1)$ and $U \sim \chi^2_\nu$ are independent, then
$T = \frac{Z}{\sqrt{U/\nu}} \sim t_\nu$, which defines the $t$-distribution. Under normality, $\sqrt{n}(\overline{X}_n-\mu)/\sigma \sim N(0,1)$ and $(n-1)S_{1n}^2/\sigma^2\sim \chi^2_{n-1}$, and the two are independent.)
\vspace{20em}


\item (CI for linear function of $\mu$) Suppose we have an i.i.d. random sample of size $n$, $X_1,...,X_n$, from $X \stackrel{iid}{\sim} N(\mu,\sigma^2)$. For constants $a,b$ with $b\neq0$, derive a $100(1-\alpha)\%$ CI for $a+b\mu$ under (i) known $\sigma^2$, (ii) unknown $\sigma^2$.

\end{enumerate}






\end{document}
